\subsection{分次模中的准素分解}

命题 5. 设 \(\Delta\) 是无挠交换群，A 是 \(\Delta\) 型分次诺特环，M 是 M 型分次 A-模，N 是 M 的分次子模，且 \(N = \bigcap_{i \in I} Q_i\) 是 N 在 M 中的准素分解。

\begin{enumerate}
\def\labelenumi{(\roman{enumi})}
\item
  令 \(\mathbf{Q}_i'\) 为 \(\mathbf{M}\) 中包含于 \(\mathbf{Q}_i\) 的最大分次子模。则 \(\mathbf{Q}_i'\) 是准素的，且 \(\mathbf{N} = \bigcap_{i \in I} \mathbf{Q}_i'\)。
\item
  如果准素分解 \(\mathbf{N} = \bigcap_{i\in I}\mathbf{Q}_i\) 是既约的，则准素分解 \(\mathbf{N} = \bigcap_{i\in I}\mathbf{Q}_i'\) 也是既约的，并且对于所有 \(i\in \mathbf{I}\)，对应于 \(\mathbf{Q}_i\) 和 \(\mathbf{Q}_i'\) 的素理想相等。
\item
  如果 \(\mathbf{Q}_i\) 对应于素理想 \(\mathfrak{p}_i\)，且该理想是 \(\operatorname{Ass}(\mathbf{M} / \mathbf{N})\) 的极小元，则 \(\mathbf{Q}_i\) 是 \(\mathbf{M}\) 的分次子模。
\end{enumerate}

  我们已经看到（第 2 号，命题 4），\(Q_{i}^{\prime}\) 关于 M 是准素的，且 \(N \subset Q_{i}^{\prime} \subset Q_{i}\)，这证明了 (i)。第 2 号的命题 4 还表明，素理想 \(p_{i}^{\prime}\) 对应于 \(Q_{i}^{\prime}\)，是包含于素理想 \(p_{i}\)（对应于 \(Q_{i}\)）中的最大分次理想。若分解 \(N = \bigcap_{i \in I} Q_{i}\) 是既约的，则有 \(p_{i} \in \text{Ass}(M/N)\) 对于所有 \(i\)（\(\S 2\)，第 3 号，命题 4），所以 \(p_{i}\) 是分次理想（第 1 号，命题 1），从而 \(p_{i}^{\prime} = p_{i}\)；于是 \(\text{Ass}(M/N) = \bigcup_{i \in I} \{p_{i}^{\prime}\}\)（\(\S 2\)，第 3 号，命题 4），这证明分解 \(N = \bigcap_{i \in I} Q_{i}^{\prime}\) 是既约的（\(\S 2\)，第 3 号，命题 4）。最后，若 \(p_{i}\) 是 \(\text{Ass}(M/N)\) 的极小元，则 \(p_{i}^{\prime} = p_{i}\)；由于 \(p_{i}\) 是分次理想（第 1 号，命题 1），从而 \(Q_{i}^{\prime} = Q_{i}\)，依据 \(\S 2\)。

% semantic-fragments: legacy-input-seam
\relax{}
